% GENERATED FILE. Edit canonical lessons, then run npm run generate:books.
% canonical-source-hash: fnv1a64:96db0b007be2e9be
% canonical-lessons: RU-C02-ya, RU-C02-ty-vy, RU-C02-vy-formality, RU-C02-menya-zovut, RU-C02-kak-vas-zovut, RU-C02-kak-cross-language, RU-C02-ochen-priyatno, RU-C02-practice, RU-C02-practice-cases, RU-C02-practice-zero-copula

\chapter{Introducing yourself}
\label{ch:russian-introductions}
\hlchaptermodality{2}

\begin{chapteropening}
\textbf{By the end of this chapter:} \emph{I can give my name in Russian, ask for someone else's, and pick \ru{ты} or \ru{вы} to match how well I know them.}

Run the whole introduction with a stranger — greet, ask \ru{Как} \ru{вас} \ru{зовут}?, answer \ru{Меня} \ru{зовут}…, close with \ru{Очень} \ru{приятно} — then switch it for one friend by changing only the greeting and \ru{вас} $\to$ \ru{тебя}.
\end{chapteropening}

\section[ya]{\ru{я} — \textquotedblleft{}I\textquotedblright{}}

\label{lesson:RU-C02-ya}



\begin{quote}
The shortest word in the chapter: one letter, one sound, and a pedigree of about six thousand years.
\end{quote}

\subsection*{You'll want to know first — \ru{я}}
\begin{quote}
\textbf{\ru{я}} — \emph{ya} — \textbf{I}
\end{quote}

That is the whole word. A single letter, pronounced \emph{ya}.

It is a \textbf{new letter} — the writing track has taught \ru{в}, \ru{р}, \ru{с}, \ru{н}, \ru{б}, \ru{д}, \ru{п}, \ru{и}, \ru{е} and \ru{т}, and hasn't reached \textbf{\ru{я}}. Read it here; you'll draw it when the track gets there.

And be careful with its shape. \textbf{\ru{я} looks like a mirrored Latin R and has nothing to do with R.} (It descends from \ru{ѧ}, the \textquotedblleft{}little yus\textquotedblright{} — not from R at all.)

It is also the \textbf{last} letter of the alphabet, 33rd of 33, and Russian has a rebuke built on that: «\textbf{\ru{Я} — \ru{последняя} \ru{буква} \ru{в} \ru{алфавите}}», \textquotedblleft{}\emph{\ru{я}} is the last letter of the alphabet.\textquotedblright{} You say it to someone who talks about themselves too much. The letter's position is the joke — don't put yourself first.

\begin{cousinweb}
Pronouns are the \textbf{least borrowable} part of a language. Nouns come and go by the thousand — \emph{sputnik}, \emph{pizza}, \emph{computer} — but nobody borrows a word for \emph{I}. So \emph{\ru{я}} has simply been inherited, without interruption, for as long as there has been an Indo-European family:

Russian \textbf{\ru{я}}, Latin \emph{ego}, German \emph{ich}, English \textbf{I}, Greek \emph{eg\'{\={o}}}.

All from PIE *\emph{eǵh\textsubscript{2}(om)}. Russian's route ran through Proto-Slavic *\emph{az\ru{ъ}}, which East Slavic reshaped — a \emph{j-} was added at the front (*\emph{az\ru{ъ}} $\to$ \emph{\ru{язъ}}) and the ending fell away — leaving the \textbf{one letter} you now say. It is the most extreme change in the family, and the reason \emph{\ru{я}} sounds nothing like \emph{ego} while being the same word.

Latin \emph{ego} is the one English borrowed back, whole and untranslated, for the part of the mind that says \emph{I}.
\end{cousinweb}

\subsection*{Your turn}
Say these aloud:

\begin{itemize}
\raggedright
  \item \textquotedblleft{}\ru{я}\textquotedblright{} — \emph{ya}
  \item the family — \textquotedblleft{}\textbf{\ru{я}} … \emph{ego} … \emph{ich} … \textbf{I}\textquotedblright{}
  \item the erosion — \textquotedblleft{}*\emph{az\ru{ъ}} … \textbf{\ru{я}}\textquotedblright{}
\end{itemize}

\begin{tcolorbox}[breakable,colback=teal!4,colframe=teal!35!black,title={Before you move on}]
Say \textquotedblleft{}I\textquotedblright{}. (\textbf{\ru{Я}} — one letter, said \emph{ya}.) What must you not mistake it for? (A \textbf{Latin R} — the shape is a coincidence.) Why has it changed so little in meaning? (\textbf{Pronouns are the least borrowable words} a language has.) Name three relatives. (\emph{Ego}, \emph{ich}, \emph{I} — and Greek \emph{eg\'{\={o}}}.) What did Russian's come from? (Proto-Slavic *\emph{az\ru{ъ}}, with a \emph{j-} added at the front and the ending lost.) Next: the two words for \textquotedblleft{}you.\textquotedblright{}
\end{tcolorbox}
\section[ty, vy]{\ru{ты}, \ru{вы} — the two words for \textquotedblleft{}you\textquotedblright{}}

\label{lesson:RU-C02-ty-vy}



\begin{quote}
Chapter 1 gave you two ways to say hello, split by formality. That split runs through the whole language — and here is the place it matters most: Russian has \textbf{two words for \textquotedblleft{}you.\textquotedblright{}}
\end{quote}

\subsection*{The two words}
\begin{quote}
\textbf{\ru{ты}} — \emph{ty} — you, one person you're familiar with. \textbf{\ru{вы}} — \emph{vy} — you, formal, \textbf{or} more than one person.
\end{quote}

Both turn on \textbf{\ru{ы}}, which is the hardest vowel in Russian for an English speaker: not the \emph{ee} of \emph{\ru{ты}}'s look-alike, but a tighter, further-back sound made with the tongue pulled towards the throat. Say English \emph{ill}, then try to say it with your tongue drawn back. The writing track hasn't reached \textbf{\ru{ы}} either.

\begin{cousinweb}
Pronouns are the \textbf{least borrowable} part of a language. Nouns come and go; \emph{I} and \emph{you} stay. So the first two line up almost too neatly:

\begin{itemize}
\raggedright
  \item \textbf{\ru{я}} (last lesson) — Latin \emph{ego}, English \textbf{I}, German \emph{ich}.
  \item \textbf{\ru{ты}} — Latin \emph{tū}, English \textbf{thou}, German \emph{du}.
  \item \textbf{\ru{вы}} — Latin \emph{vōs}, with no English or German survivor to set beside it.
\end{itemize}

The first two lines are the same word four times over. \textbf{The third is only half a set}: \emph{\ru{вы}} and \emph{vōs} both continue PIE *\emph{wos}, but English \emph{you} and German \emph{ihr} come from the \textbf{other} stem in that paradigm, *\emph{yūs}. Cousins, not twins — worth marking rather than smoothing over.

\textbf{\ru{ты}} and \emph{tū} and \emph{thou} are the same word, barely touched in four thousand years. \emph{Thou} fell out of standard English through the 1600s (it survives in Quaker usage, northern dialect and liturgy), leaving \emph{you} for everyone — which is why the \ru{ты}/\ru{вы} choice feels foreign to English speakers. English \textbf{had} it.
\end{cousinweb}

\subsection*{Your turn}
Say these aloud:

\begin{itemize}
\raggedright
  \item \textquotedblleft{}\ru{я}\textquotedblright{} — \emph{ya}
  \item \textquotedblleft{}\ru{ты}\textquotedblright{} and \textquotedblleft{}\ru{вы}\textquotedblright{} — hear the difference
  \item the family — \textquotedblleft{}\textbf{\ru{ты}} … \emph{tū} … \textbf{thou}\textquotedblright{}
\end{itemize}

\begin{tcolorbox}[breakable,colback=teal!4,colframe=teal!35!black,title={Before you move on}]
Say \textquotedblleft{}I\textquotedblright{}. (\textbf{\ru{Я}}.) What are the two words for \textquotedblleft{}you\textquotedblright{}, and what splits them? (\textbf{\ru{Ты}} informal, \textbf{\ru{вы}} formal — or plural.) Why is \emph{\ru{ты}} familiar to an English speaker's ear? (It's the \textbf{same word as \emph{thou}}, which left standard English through the 1600s.) Next: why plural \emph{\ru{вы}} became polite.
\end{tcolorbox}
\section[vy]{Why plural \emph{\ru{вы}} became polite}

\label{lesson:RU-C02-vy-formality}



\begin{quote}
Give the pair: \emph{\ru{ты}} for one familiar person; \emph{\ru{вы}} for more than one person — or for one person addressed respectfully.
\end{quote}

\begin{grammarlens}[title={Grammar Lens: one social job, several grammatical routes}]
Russian \textbf{\ru{вы}} is a plural used for \textbf{one person} as a mark of respect. French does the same job with second-person plural \textbf{vous}. Chapter 1's \emph{\ru{здравствуйте}} already carried the matching polite-plural ending \emph{-\ru{те}}.

Other languages reached respect differently:

\begin{itemize}
\raggedright
  \item Russian \textbf{\ru{вы}} — the second-person plural.
  \item French \textbf{vous} — the second-person plural, the very same route.
  \item German \textbf{Sie} — the \emph{third}-person plural, not the ordinary \emph{ihr}.
  \item Spanish \textbf{usted} — a worn-down title, \emph{vuestra merced}, \textquotedblleft{}your grace.\textquotedblright{}
\end{itemize}

Russian's polite \emph{\ru{вы}} is generally treated as an eighteenth-century import from French and German social usage, not a pattern Russian invented alone.
\end{grammarlens}

\begin{culture}
Use \emph{\ru{вы}} with a stranger, an older person, or in a formal relationship. Use \emph{\ru{ты}} with a friend, child, or family member. When unsure with an adult, begin with \textbf{\ru{вы}}; too much respect is safer than unwanted familiarity.
\end{culture}

\subsection*{Your turn}
Say these aloud:

\begin{itemize}
\raggedright
  \item one close friend — \emph{\ru{ты}}
  \item a new professor — \emph{\ru{вы}}
  \item Russian and French use which route — second-person plural
  \item Spanish \emph{usted} came from what — \textquotedblleft{}your grace\textquotedblright{}
\end{itemize}

\begin{tcolorbox}[breakable,colback=teal!4,colframe=teal!35!black,title={Before you move on}]
Why can \emph{\ru{вы}} address one person? (A plural became respectful.) Which language shares that exact route? (French \emph{vous}.) Which two take different routes? (German \emph{Sie} uses third-person plural; Spanish \emph{usted} came from a title.) Which Russian form is safest when you are unsure? (\emph{\ru{Вы}}.)
\end{tcolorbox}
\section[menyá zovút]{\ru{меня} \ru{зовут}… — \textquotedblleft{}my name is…\textquotedblright{}}

\label{lesson:RU-C02-menya-zovut}



\begin{quote}
Many languages say \textquotedblleft{}my name is\textquotedblright{} with a word for \emph{my} and a word for \emph{name}. Russian says something quite different, and it is worth taking apart slowly.
\end{quote}

\subsection*{The exchange}
\begin{quote}
\textbf{\ru{Меня} \ru{зовут} \ru{Анна}.} — \emph{Menyá zovút Anna.} — \textquotedblleft{}My name is Anna.\textquotedblright{}
\end{quote}

Now the literal reading:

\textbf{\ru{меня}} — \emph{menya} — \textbf{me}, the object form. \textbf{\ru{зовут}} — \emph{zovut} — \textbf{they call}.

So: \textbf{\textquotedblleft{}{[}they{]} call me Anna.\textquotedblright{}}

There is \textbf{no word for \textquotedblleft{}my\textquotedblright{}} and \textbf{no word for \textquotedblleft{}name\textquotedblright{}} in the sentence. Russian reports what people call you, not what you possess.

\begin{grammarlens}[title={Grammar Lens: the invisible \textquotedblleft{}they\textquotedblright{}}]
\emph{\ru{Зовут}} is the \textbf{they} form of \emph{\ru{звать}}, \textquotedblleft{}to call.\textquotedblright{} But there is no \emph{\ru{они}} (\textquotedblleft{}they\textquotedblright{}) anywhere — and there isn't meant to be. Russian uses a bare plural verb to mean \textquotedblleft{}people in general,\textquotedblright{} the way English says \emph{\textquotedblleft{}they say it'll rain\textquotedblright{}} without anyone in mind.

English has the same trick and doesn't notice it:

\begin{quote}
\emph{They call me Anna.} — Who are \emph{they}? Nobody. Everybody.
\end{quote}
\end{grammarlens}

\begin{grammarlens}[title={Grammar Lens: \ru{меня} is not \ru{я}}]
This is the important part, and this book's first taste of \textbf{case}.

You learned \textbf{\ru{я}} for \textquotedblleft{}I\textquotedblright{} last lesson. But here it is \textbf{\ru{меня}} — because \emph{I} is now the \textbf{object} of the calling, not the one doing it. Russian changes the \textbf{shape of the word} to mark that job:

\textbf{\ru{я}} is \emph{I}, the one acting; \textbf{\ru{меня}} is \emph{me}, the one acted on.

English does this with about six \textbf{pairs} — \emph{I/me}, \emph{he/him}, \emph{she/her}, \emph{we/us}, \emph{they/them}, \emph{who/whom}, plus the archaic \emph{thou/thee}. Russian does it with \textbf{every noun and pronoun in the language}. That system is called \textbf{case}, and it is the main thing standing between you and Russian.

You do not need to learn it yet. You need to notice that \emph{\ru{меня}} is a \textbf{different shape of \ru{я}}, and that the shape carries the meaning.
\end{grammarlens}

\subsection*{Your turn}
Say these aloud:

\begin{itemize}
\raggedright
  \item \textquotedblleft{}\ru{Меня} \ru{зовут}…\textquotedblright{} and your own name
  \item the literal — \textquotedblleft{}\textbf{they call me} …\textquotedblright{}
  \item the pair — \textquotedblleft{}\textbf{\ru{я}} … \textbf{\ru{меня}}\textquotedblright{}
  \item the English echo — \textquotedblleft{}\emph{they say it'll rain}\textquotedblright{}
\end{itemize}

\begin{tcolorbox}[breakable,colback=teal!4,colframe=teal!35!black,title={Before you move on}]
Say \textquotedblleft{}my name is Anna.\textquotedblright{} (\emph{\ru{Меня} \ru{зовут} \ru{Анна}}.) What does it literally mean? (\textquotedblleft{}\textbf{They call me} Anna.\textquotedblright{}) Which two words is it missing? (Any word for \textbf{my}, and any word for \textbf{name}.) Who are \textquotedblleft{}they\textquotedblright{}? (\textbf{Nobody} — a bare plural verb meaning \textquotedblleft{}people in general,\textquotedblright{} like English \emph{they say}.) Why \emph{\ru{меня}} and not \emph{\ru{я}}? (Because it's the \textbf{object} — Russian changes a word's \textbf{shape} for its job. That's \textbf{case}.) Next: asking someone else's name.
\end{tcolorbox}
\section[kak vas zovút]{\ru{как} \ru{вас} \ru{зовут}? — \textquotedblleft{}what's your name?\textquotedblright{}}

\label{lesson:RU-C02-kak-vas-zovut}



\begin{quote}
You can say your own name. Now ask for someone else's — and notice that Russian asks a different question from English.
\end{quote}

\subsection*{The exchange}
\begin{quote}
\textbf{\ru{Как} \ru{вас} \ru{зовут}?} — \emph{Kak vas zovút?} — formal, or to more than one person. \textbf{\ru{Как} \ru{тебя} \ru{зовут}?} — \emph{Kak tebyá zovút?} — informal, to one person you know.
\end{quote}

Taken apart:

\textbf{\ru{как}} — \emph{kak} — \textbf{how}. \textbf{\ru{вас}} or \textbf{\ru{тебя}} — \emph{vas} / \emph{tebya} — \textbf{you}, in the object form. \textbf{\ru{зовут}} — \emph{zovut} — they call.

So: \textbf{\textquotedblleft{}How do they call you?\textquotedblright{}}

\begin{grammarlens}[title={Grammar Lens: the object forms again}]
Last lesson's \emph{\ru{меня}} has partners, and they follow the same logic — the person being called is the \textbf{object}, so the word changes shape:

\begin{itemize}
\raggedright
  \item \textbf{\ru{я}} becomes \textbf{\ru{меня}} — me.
  \item \textbf{\ru{ты}} becomes \textbf{\ru{тебя}} — you, informal.
  \item \textbf{\ru{вы}} becomes \textbf{\ru{вас}} — you, formal or plural.
\end{itemize}

You now have a matched pair: \textbf{\ru{Как} \ru{вас} \ru{зовут}?} — \textbf{\ru{Меня} \ru{зовут}…}
\end{grammarlens}

\subsection*{Your turn}
Say these aloud:

\begin{itemize}
\raggedright
  \item \textquotedblleft{}\ru{Как} \ru{вас} \ru{зовут}?\textquotedblright{} — formal
  \item \textquotedblleft{}\ru{Как} \ru{тебя} \ru{зовут}?\textquotedblright{} — informal
  \item the whole exchange — \textquotedblleft{}\ru{Как} \ru{вас} \ru{зовут}?\textquotedblright{} · \textquotedblleft{}\ru{Меня} \ru{зовут}…\textquotedblright{}
\end{itemize}

\begin{tcolorbox}[breakable,colback=teal!4,colframe=teal!35!black,title={Before you move on}]
Ask a stranger's name. (\emph{\ru{Как} \ru{вас} \ru{зовут}?}) And a friend's? (\emph{\ru{Как} \ru{тебя} \ru{зовут}?}) What does \emph{\ru{как}} mean, and why does it matter? (\textbf{How} — Russian asks \emph{how they call you}, about an \textbf{action}, where English asks \emph{what}, about a \textbf{possession}.) Give the three object forms. (\emph{\ru{Меня}, \ru{тебя}, \ru{вас}}.) Next: compare how several languages frame the same question.
\end{tcolorbox}
\section[kak / comment / como]{Why Russian asks \emph{how}, not \emph{what}}

\label{lesson:RU-C02-kak-cross-language}



\begin{quote}
In \emph{\ru{Как} \ru{вас} \ru{зовут}?}, which word means \textquotedblleft{}how\textquotedblright{}? (\emph{\ru{Как}}.)
\end{quote}

\begin{grammarlens}[title={Grammar Lens: one question, two ways to organize it}]
English asks \textbf{what} your name is, treating the name as a thing you have. Russian asks \textbf{how} people call you, treating naming as something people do.

English is the odd one out in this small comparison:

\begin{itemize}
\raggedright
  \item Russian asks \emph{\ru{Как} \ru{вас} \ru{зовут}?} — \textbf{how}.
  \item French asks \emph{Comment vous appelez-vous ?} — \textbf{how}.
  \item Spanish asks \emph{¿Cómo se llama?} — \textbf{how}.
  \item English asks \emph{\textbf{What}} is your name? — \textbf{what}.
\end{itemize}

Three languages organize the question around an \textbf{action}; English organizes it around a \textbf{possession}. The answer is the same social fact, but the route to it differs.
\end{grammarlens}

\subsection*{Your turn}
Say these aloud:

\begin{itemize}
\raggedright
  \item the three action words — \emph{\ru{как}}, \emph{comment}, \emph{cómo}
  \item Russian literally asks — \textquotedblleft{}how do they call you?\textquotedblright{}
  \item English asks about — a thing you have
\end{itemize}

\begin{tcolorbox}[breakable,colback=teal!4,colframe=teal!35!black,title={Before you move on}]
Which three languages ask \textquotedblleft{}how\textquotedblright{}? (Russian, French, Spanish.) Which asks \textquotedblleft{}what\textquotedblright{}? (English.) What is the underlying contrast? (Naming as an action versus a name as a possession.)
\end{tcolorbox}
\section[óchen priyátno]{\ru{очень} \ru{приятно} — \textquotedblleft{}pleased to meet you\textquotedblright{}}

\label{lesson:RU-C02-ochen-priyatno}



\begin{quote}
The reply that closes an introduction — and the best etymology in the chapter.
\end{quote}

\subsection*{The two words}
\begin{quote}
\textbf{\ru{Очень} \ru{приятно}.} — \emph{Óchen priyátno.} — \textquotedblleft{}Pleased to meet you.\textquotedblright{}
\end{quote}

Literally: \textbf{\textquotedblleft{}very pleasant.\textquotedblright{}}

\textbf{\ru{очень}} — \emph{ochen} — very. \textbf{\ru{приятно}} — \emph{priyatno} — pleasant, agreeable.

There is no \textquotedblleft{}to meet you\textquotedblright{} and no \textquotedblleft{}I am.\textquotedblright{} Russian states the \textbf{quality of the moment} and leaves the rest understood — the same economy as \emph{\ru{меня} \ru{зовут}}, which left out \textquotedblleft{}my\textquotedblright{} and \textquotedblleft{}name.\textquotedblright{}

Both \textbf{\ru{ч}} (\emph{ch}) and \textbf{\ru{я}} are letters the writing track hasn't reached. Read them; you'll draw them later.

\begin{cousinweb}
\emph{\ru{Приятно}} comes from Slavic \emph{\textbf{prijati}}, \textquotedblleft{}to favour, to be well-disposed toward\textquotedblright{} — which gives Russian \textbf{\ru{приятель}} (\emph{priyátel'}), a friend.

That root goes back to PIE *\emph{preyH-}, \textquotedblleft{}to love, to please.\textquotedblright{} And it went into Germanic too:

\begin{itemize}
\raggedright
  \item Russian \textbf{\ru{приятно}}, pleasant, and \textbf{\ru{приятель}}, a friend.
  \item English \textbf{friend} — literally \textquotedblleft{}one who is loving,\textquotedblright{} an old participle of a verb meaning \emph{to love}.
  \item German \emph{Freund}.
\end{itemize}

So \emph{\ru{приятно}} and \emph{friend} are the same root, arriving from opposite ends of Europe.

\textbf{Free} almost certainly belongs to the family too. The usual account: the word first meant \textquotedblleft{}belonging to the beloved household\textquotedblright{} — the members of a family as opposed to the slaves in it — and only later meant \textquotedblleft{}not in bondage.\textquotedblright{} \emph{Free} and \emph{friend} start in the same place.

Saying \emph{\ru{очень} \ru{приятно}} is, at the root, calling the meeting \textbf{friendly}.
\end{cousinweb}

\subsection*{Your turn}
Say these aloud:

\begin{itemize}
\raggedright
  \item \textquotedblleft{}\ru{Очень} \ru{приятно}\textquotedblright{}
  \item the whole exchange — \textquotedblleft{}\ru{Как} \ru{вас} \ru{зовут}?\textquotedblright{} · \textquotedblleft{}\ru{Меня} \ru{зовут}…\textquotedblright{} · \textquotedblleft{}\ru{Очень} \ru{приятно}\textquotedblright{}
  \item the family — \textquotedblleft{}\textbf{\ru{прия}}\ru{тно} … \textbf{\ru{прия}}\ru{тель} … \textbf{friend}\textquotedblright{}
\end{itemize}

\begin{tcolorbox}[breakable,colback=teal!4,colframe=teal!35!black,title={Before you move on}]
Say \textquotedblleft{}pleased to meet you.\textquotedblright{} (\emph{\ru{Очень} \ru{приятно}}.) What does it literally mean? (\textquotedblleft{}\textbf{Very pleasant}\textquotedblright{} — no \textquotedblleft{}I am,\textquotedblright{} no \textquotedblleft{}to meet you.\textquotedblright{}) What is \emph{\ru{приятно}} related to? (Slavic \emph{prijati} \textquotedblleft{}to favour\textquotedblright{} $\to$ \textbf{\ru{приятель}} \textquotedblleft{}friend\textquotedblright{} — and, from the same PIE root *\emph{preyH-}, English \textbf{friend} and German \emph{Freund}.) And \emph{free}? (Most likely the \textbf{same family} — originally \textquotedblleft{}belonging to the beloved household,\textquotedblright{} i.e. not a slave.) Next: putting the chapter together.
\end{tcolorbox}
\section[Practice]{Chapter 2 practice — run the exchange}

\label{lesson:RU-C02-practice}



\begin{quote}
Give the three words the chapter opened with: \emph{I}, familiar \emph{you}, formal \emph{you}. (\emph{\ru{Я}}, \emph{\ru{ты}}, \emph{\ru{вы}}.) Everything below is built from those.
\end{quote}

\subsection*{The exchange}
Two strangers meet. Say both voices aloud, one after the other:

\begin{quote}
— \textbf{\ru{Здравствуйте}!} — \emph{Zdrávstvuyte!} — \textbf{\ru{Как} \ru{вас} \ru{зовут}?} — \emph{Kak vas zovút?} — \textbf{\ru{Меня} \ru{зовут} \ru{Анна}. \ru{А} \ru{вас}?} — \emph{Menyá zovút Anna. A vas?} — \textbf{\ru{Меня} \ru{зовут} \ru{Иван}. \ru{Очень} \ru{приятно}.} — \emph{Menyá zovút Iván. Óchen priyátno.}
\end{quote}

Everything in it is formal, because \emph{\ru{вы}} is the safe opening with an adult you do not know.

\begin{grammarlens}[title={Grammar Lens: what moves when the exchange goes informal}]
With one friend, exactly two pieces change:

\begin{itemize}
\raggedright
  \item the greeting — \textbf{\ru{Здравствуйте}} becomes \textbf{\ru{Привет}};
  \item the object pronoun — \textbf{\ru{вас}} becomes \textbf{\ru{тебя}}.
\end{itemize}

So \emph{\ru{Как} \ru{вас} \ru{зовут}?} becomes \emph{\ru{Как} \ru{тебя} \ru{зовут}?} The verb \textbf{\ru{зовут}} does not move, and \textbf{\ru{меня}} in the answer does not move either — you are still the one being called, whoever is asking.
\end{grammarlens}

\begin{grammarlens}[title={Grammar Lens: why the shapes change at all}]
\emph{\ru{Я}} becomes \emph{\ru{меня}} because the calling lands \textbf{on} you rather than coming \textbf{from} you; \emph{\ru{ты}} and \emph{\ru{вы}} become \emph{\ru{тебя}} and \emph{\ru{вас}} for the same reason.

And \emph{\ru{зовут}} is a plural verb with nobody behind it — \textquotedblleft{}\textbf{they} call me,\textquotedblright{} where \emph{they} is no one in particular. That is why the Russian question asks \textbf{how} people call you, not \textbf{what} your name is: naming is something people do, not something you own.
\end{grammarlens}

\subsection*{Your turn}
Say these aloud:

\begin{itemize}
\raggedright
  \item the full formal exchange, both voices
  \item the same exchange to one friend
  \item your own answer — \textquotedblleft{}\ru{Меня} \ru{зовут}…\textquotedblright{} and your name
\end{itemize}

\begin{tcolorbox}[breakable,colback=teal!4,colframe=teal!35!black,title={Before you move on}]
Run the formal exchange from memory, then switch it for one friend. Which two pieces changed? (The greeting, and \emph{\ru{вас} $\to$ \ru{тебя}}.) Which word never moved? (\emph{\ru{Зовут}}.) Which form do you choose for an adult you have just met, and why? (\emph{\ru{Вы}} — a plural used as respect, and the safer default.) Who is doing the calling in \emph{\ru{меня} \ru{зовут}}? (Nobody in particular — a bare plural verb.) And why is the question \emph{\ru{как}} rather than \emph{\ru{что}}? (Russian asks \textbf{how} people call you.)
\end{tcolorbox}
\section[Practice]{Chapter 2 practice — the two person shapes}

\label{lesson:RU-C02-practice-cases}



\begin{quote}
Russian changes a pronoun's shape when the action lands on that person. English still shows the same idea in \emph{I $\to$ me}.
\end{quote}

\begin{grammarlens}[title={Grammar lens — cover and retrieve}]
Cover the right column, then give the acted-on form:

\noindent
\begin{tabularx}{\linewidth}{@{}>{\raggedright\arraybackslash}X>{\raggedright\arraybackslash}X>{\raggedright\arraybackslash}X@{}}
\toprule
\textbf{doing the action} & \textbf{action lands here} & \textbf{meaning} \\
\midrule
\textbf{\ru{я}} & \textbf{\ru{меня}} & I / me \\
\textbf{\ru{ты}} & \textbf{\ru{тебя}} & you / you, informal \\
\textbf{\ru{вы}} & \textbf{\ru{вас}} & you / you, formal or plural \\
\bottomrule
\end{tabularx}

Russian has six cases, but none of that system is due yet. Learn these three pairs because the naming exchange needs them.
\end{grammarlens}

\subsection*{How to answer — each shape in the frame}
\begin{itemize}
\raggedright
  \item \textbf{\ru{Меня}} \ru{зовут}… — people call \textbf{me}…
  \item \ru{Как} \textbf{\ru{тебя}} \ru{зовут}? — how do people call \textbf{you}, one friend?
  \item \ru{Как} \textbf{\ru{вас}} \ru{зовут}? — how do people call \textbf{you}, formally?
\end{itemize}

The verb stays \textbf{\ru{зовут}}; the person changes shape.

\subsection*{Your turn}
Say these aloud:

\begin{itemize}
\raggedright
  \item \emph{\ru{я} $\to$ \ru{меня}}
  \item \emph{\ru{ты} $\to$ \ru{тебя}}
  \item \emph{\ru{вы} $\to$ \ru{вас}}
  \item one verb, three object forms
\end{itemize}

\begin{tcolorbox}[breakable,colback=teal!4,colframe=teal!35!black,title={Before you move on}]
Give the three pairs. Which shape answers \textquotedblleft{}who receives the calling\textquotedblright{}? (The object form.) Do you need the six-case table yet? (No — only these earned pairs.)
\end{tcolorbox}
\section[Practice]{Chapter 2 practice — what Russian leaves out}

\label{lesson:RU-C02-practice-zero-copula}



\begin{quote}
Russian often says less than the English translation. Name only the words that are genuinely absent; do not erase a verb that is really there.
\end{quote}

\begin{grammarlens}[title={Grammar lens — three phrases, three English expansions}]
\begin{itemize}
\raggedright
  \item \textbf{\ru{Меня} \ru{зовут} \ru{Анна}} reads literally \textquotedblleft{}they call me Anna\textquotedblright{}; English adds \emph{my} and \emph{name}.
  \item \textbf{\ru{Очень} \ru{приятно}} reads literally \textquotedblleft{}very pleasant\textquotedblright{}; English adds \emph{it is} and \emph{to meet you}.
  \item \textbf{\ru{Как} \ru{вас} \ru{зовут}?} reads literally \textquotedblleft{}how do they call you?\textquotedblright{}; English adds \emph{what} and \emph{your}.
\end{itemize}

Only the middle sentence has \textbf{no verb at all}. In present-tense descriptions, Russian normally drops the linking verb: \emph{\ru{Очень} \ru{приятно}} needs no spoken \textquotedblleft{}is.\textquotedblright{}

The other two sentences do contain \textbf{\ru{зовут}}, \textquotedblleft{}they call.\textquotedblright{} They lack \textquotedblleft{}is\textquotedblright{} because they are not \textquotedblleft{}X is Y\textquotedblright{} sentences in Russian; English simply translates them that way.

Russian also has a present-tense \textbf{\ru{есть}}. Chapter 1 already hid it inside \textbf{\ru{нет}} = \emph{\ru{не} + \ru{есть}}, \textquotedblleft{}not-is.\textquotedblright{} The next chapter will reuse it. So the careful rule is not \textquotedblleft{}Russian has no \emph{is}\textquotedblright{}; it is \textquotedblleft{}Russian usually omits the linking verb in a present-tense description.\textquotedblright{}
\end{grammarlens}

\subsection*{Your turn}
Say these aloud:

\begin{itemize}
\raggedright
  \item the one verb-free phrase — \emph{\ru{Очень} \ru{приятно}}
  \item the verb inside the naming pair — \emph{\ru{зовут}}
  \item the verb hidden inside \emph{\ru{нет}} — \emph{\ru{есть}}
\end{itemize}

\begin{tcolorbox}[breakable,colback=teal!4,colframe=teal!35!black,title={Before you move on}]
Which sentence demonstrates the zero copula? (\emph{\ru{Очень} \ru{приятно}}.) Why do the naming sentences not count? (They contain \emph{\ru{зовут}}.) Does Russian lack a present form of \textquotedblleft{}be\textquotedblright{} entirely? (No: \emph{\ru{есть}} exists; the linking use is normally omitted here.)
\end{tcolorbox}
